\documentclass{amsart}
% For dealing with references we use the comment environment
\usepackage{verbatim}
\newenvironment{reference}{\comment}{\endcomment}
%\newenvironment{reference}{}{}
\newenvironment{slogan}{\comment}{\endcomment}
\newenvironment{history}{\comment}{\endcomment}

% For commutative diagrams we use Xy-pic
\usepackage[all]{xy}

% We use 2cell for 2-commutative diagrams.
\xyoption{2cell}
\UseAllTwocells

% We use multicol for the list of chapters between chapters
\usepackage{multicol}

% This is generally recommended for better output
\usepackage{lmodern}
\usepackage[T1]{fontenc}

% For cross-file-references

% Package for hypertext links:
\usepackage{hyperref}

% For any local file, say "hello.tex" you want to link to please

% Theorem environments.
%
\theoremstyle{plain}
\newtheorem{theorem}[subsection]{Theorem}
\newtheorem{proposition}[subsection]{Proposition}
\newtheorem{lemma}[subsection]{Lemma}

\theoremstyle{definition}
\newtheorem{definition}[subsection]{Definition}
\newtheorem{example}[subsection]{Example}
\newtheorem{exercise}[subsection]{Exercise}
\newtheorem{situation}[subsection]{Situation}

\theoremstyle{remark}
\newtheorem{remark}[subsection]{Remark}
\newtheorem{remarks}[subsection]{Remarks}

\numberwithin{equation}{subsection}

% Macros
%
\def\lim{\mathop{\mathrm{lim}}\nolimits}
\def\colim{\mathop{\mathrm{colim}}\nolimits}
\def\Spec{\mathop{\mathrm{Spec}}}
\def\Hom{\mathop{\mathrm{Hom}}\nolimits}
\def\Ext{\mathop{\mathrm{Ext}}\nolimits}
\def\SheafHom{\mathop{\mathcal{H}\!\mathit{om}}\nolimits}
\def\SheafExt{\mathop{\mathcal{E}\!\mathit{xt}}\nolimits}
\def\Sch{\mathit{Sch}}
\def\Mor{\mathop{\mathrm{Mor}}\nolimits}
\def\Ob{\mathop{\mathrm{Ob}}\nolimits}
\def\Sh{\mathop{\mathit{Sh}}\nolimits}
\def\NL{\mathop{N\!L}\nolimits}
\def\CH{\mathop{\mathrm{CH}}\nolimits}
\def\proetale{{pro\text{-}\acute{e}tale}}
\def\etale{{\acute{e}tale}}
\def\QCoh{\mathit{QCoh}}
\def\Ker{\mathop{\mathrm{Ker}}}
\def\Im{\mathop{\mathrm{Im}}}
\def\Coker{\mathop{\mathrm{Coker}}}
\def\Coim{\mathop{\mathrm{Coim}}}

% Boxtimes
%
\DeclareMathSymbol{\boxtimes}{\mathbin}{AMSa}{"02}

%
% Macros for moduli stacks/spaces
%
\def\QCohstack{\mathcal{QC}\!\mathit{oh}}
\def\Cohstack{\mathcal{C}\!\mathit{oh}}
\def\Spacesstack{\mathcal{S}\!\mathit{paces}}
\def\Quotfunctor{\mathrm{Quot}}
\def\Hilbfunctor{\mathrm{Hilb}}
\def\Curvesstack{\mathcal{C}\!\mathit{urves}}
\def\Polarizedstack{\mathcal{P}\!\mathit{olarized}}
\def\Complexesstack{\mathcal{C}\!\mathit{omplexes}}
% \Pic is the operator that assigns to X its picard group, usage \Pic(X)
% \Picardstack_{X/B} denotes the Picard stack of X over B
% \Picardfunctor_{X/B} denotes the Picard functor of X over B
\def\Pic{\mathop{\mathrm{Pic}}\nolimits}
\def\Picardstack{\mathcal{P}\!\mathit{ic}}
\def\Picardfunctor{\mathrm{Pic}}
\def\Deformationcategory{\mathcal{D}\!\mathit{ef}}

\setlength{\emergencystretch}{3em}
\begin{document}
\title[Stacks Project, Tag 053U]{483 --- Closed fibre points interpolating a specialization are dense for locally finite-type morphisms}
\maketitle
\noindent Copyright (C) 2005--2025 Johan de Jong. Stacks Project authors. Source: \href{https://stacks.math.columbia.edu/tag/053U}{Tag 053U}.

\noindent Editorial difficulty note (not author text). Difficulty H2: A valuation-ring reduction must preserve the original generic fibre and the closedness of the special point. Dimension induction then chooses a principal slice avoiding arbitrary closed subsets and all special-fibre components, using the vertical-divisor criterion to ensure a horizontal component through the special point..

\noindent Editorial provenance note (not author text). History: excerpt from revision \texttt{a0444\allowbreak 6e57e\allowbreak c1fbc\allowbreak 252a8\allowbreak 71afc\allowbreak ec775\allowbreak 2fb28\allowbreak 07b14}, more-morphisms.tex, lines 14816--14934. Reference presentation was adapted; original-author context and core support blocks were retained or added. The mathematical wording is unchanged. Licensed under GNU FDL 1.2 or later; no invariant sections or cover texts. See ../licenses/COPYING. Context blocks reproduce author definitions and supporting results; they are not additional counted problems.

\begin{definition}
\label{schemes-definition-reduced-induced-scheme}
Let $X$ be a scheme. Let $Z \subset X$ be a closed subset.
A {\it scheme structure on $Z$} is given by a closed subscheme $Z'$ of
$X$ whose underlying set is equal to $Z$. We often say
``let $(Z, \mathcal{O}_Z)$ be a scheme structure on $Z$'' to
indicate this. The {\it reduced induced scheme structure}
on $Z$ is the one constructed in Lemma \href{https://stacks.math.columbia.edu/tag/01J3}{lemma reduced closed subscheme (Tag 01J3)}.
The {\it reduction $X_{red}$ of $X$} is the reduced induced scheme
structure on $X$ itself.
\end{definition}

\begin{lemma}
\label{lemma-closed-point-nearby-fibre}
Let $f : X \to S$ be a morphism of schemes.
Let $x \leadsto x'$ be a specialization of points in $X$.
Set $s = f(x)$ and $s' = f(x')$.
Assume
\begin{enumerate}
\item $x'$ is a closed point of $X_{s'}$, and
\item $f$ is locally of finite type.
\end{enumerate}
Then the set
$$
\{x_1 \in X
\text{ such that }
f(x_1) = s
\text{ and }
x_1\text{ is closed in }X_s
\text{ and }
x \leadsto x_1 \leadsto x'
\}
$$
is dense in the closure of $x$ in $X_s$.
\end{lemma}

\begin{proof}
We apply
Schemes, Lemma \href{https://stacks.math.columbia.edu/tag/01J8}{lemma points specialize (Tag 01J8)}
to the specialization $x \leadsto x'$.
This produces a morphism $\varphi : \Spec(B) \to X$
where $B$ is a valuation ring such that $\varphi$ maps the
generic point to $x$ and the closed point to $x'$. We may also
assume that $\kappa(x)$ is the fraction field of $B$.
Let $A = B \cap \kappa(s)$. Note that this is a valuation ring (see
Algebra, Lemma \href{https://stacks.math.columbia.edu/tag/052L}{lemma valuation ring cap field (Tag 052L)})
which dominates the image of $\mathcal{O}_{S, s'} \to \kappa(s)$.
Consider the commutative diagram
$$
\xymatrix{
\Spec(B) \ar[rd] \ar[r] &
X_A \ar[d] \ar[r] & X \ar[d] \\
& \Spec(A) \ar[r] & S
}
$$
The generic (resp.\ closed) point of $B$ maps to a point $x_A$
(resp.\ $x'_A$) of $X_A$ lying over the generic (resp.\ closed)
point of $\Spec(A)$. Note that $x'_A$ is a closed point
of the special fibre of $X_A$ by
Morphisms,
Lemma \href{https://stacks.math.columbia.edu/tag/053M}{lemma base change closed point fibre locally finite type (Tag 053M)}.
Note that the generic fibre of $X_A \to \Spec(A)$ is isomorphic
to $X_s$. Thus we have reduced the lemma to the case where $S$ is
the spectrum of a valuation ring, $s = \eta \in S$ is the generic point, and
$s' \in S$ is the closed point.

\medskip\noindent
We will prove the lemma by induction on $\dim_x(X_\eta)$.
If $\dim_x(X_\eta) = 0$, then there are no other points of $X_\eta$
specializing to $x$ and $x$ is closed in its fibre, see
Morphisms, Lemma \href{https://stacks.math.columbia.edu/tag/01TH}{lemma quasi finite at point characterize (Tag 01TH)},
and the result holds. Assume $\dim_x(X_\eta) > 0$.

\medskip\noindent
Let $X' \subset X$ be the reduced induced scheme structure on
the irreducible closed subscheme $\overline{\{x\}}$ of $X$, see
Schemes, Definition \ref{schemes-definition-reduced-induced-scheme}.
To prove the lemma we may replace $X$ by $X'$ as this only decreases
$\dim_x(X_\eta)$. Hence we may also assume that $X$ is an integral scheme
and that $x$ is its generic point. In addition, we may replace $X$ by an
affine neighbourhood of $x'$. Thus we have $X = \Spec(B)$ where
$A \subset B$ is a finite type extension of domains. Note that in
this case $\dim_x(X_\eta) = \dim(X_\eta) = \dim(X_{s'})$, and that in fact
$X_{s'}$ is equidimensional, see
Algebra,
Lemma \href{https://stacks.math.columbia.edu/tag/00QK}{lemma finite type domain over valuation ring dim fibres (Tag 00QK)}.

\medskip\noindent
Let $W \subset X_\eta$ be a proper closed subset (this is the
subset we want to ``avoid''). As $X_s$ is of finite type over a field
we see that $W$ has finitely many irreducible components
$W = W_1 \cup \ldots \cup W_n$. Let
$\mathfrak q_j \subset B$, $j = 1, \ldots, r$
be the corresponding prime ideals. Let $\mathfrak q \subset B$
be the maximal ideal corresponding to the point $x'$.
Let $\mathfrak p_1, \ldots, \mathfrak p_s \subset B$ be the
minimal primes lying over $\mathfrak m_AB$. There are finitely
many as these correspond to the irreducible components of the
Noetherian scheme $X_{s'}$. Moreover, each of these irreducible
components has dimension $> 0$ (see above) hence we see that
$\mathfrak p_i \not = \mathfrak q$ for all $i$.
Now, pick an element $g \in \mathfrak q$ such that
$g \not \in \mathfrak q_j$ for all $j$ and $g \not \in \mathfrak p_i$
for all $i$, see
Algebra, Lemma \href{https://stacks.math.columbia.edu/tag/00DS}{lemma silly (Tag 00DS)}.
Denote $Z \subset X$ the locally principal closed subscheme defined by $g$.
Let $Z_\eta = Z_{1, \eta} \cup \ldots \cup Z_{n, \eta}$, $n \geq 0$
be the decomposition of the generic fibre of $Z$ into irreducible
components (finitely many as the generic fibre is Noetherian).
Denote $Z_i \subset X$ the closure of $Z_{i, \eta}$.
After replacing $X$ by a smaller affine neighbourhood
we may assume that $x' \in Z_i$ for each $i = 1, \ldots, n$.
By construction $Z \cap X_{s'}$ does not contain any irreducible
component of $X_{s'}$. Hence by
Lemma \href{https://stacks.math.columbia.edu/tag/053R}{lemma locally principal vertical (Tag 053R)}
we conclude that $Z_\eta \not = \emptyset$! In other words
$n \geq 1$. Letting $x_1 \in Z_1$ be the generic point we see
that $x_1 \leadsto x'$ and $f(x_1) = \eta$.
Also, by construction $Z_{1, \eta} \cap W_j \subset W_j$
is a proper closed subset. Hence every irreducible component of
$Z_{1, \eta} \cap W_j$ has codimension $\geq 2$ in $X_\eta$
whereas $\text{codim}(Z_{1, \eta}, X_\eta) = 1$ by
Algebra, Lemma \href{https://stacks.math.columbia.edu/tag/00KV}{lemma minimal over 1 (Tag 00KV)}.
Thus $W \cap Z_{1, \eta}$ is a proper closed subset.
At this point we see that the induction hypothesis applies to
$Z_1 \to S$ and the specialization $x_1 \leadsto x'$.
This produces a closed point $x_2$ of $Z_{1, \eta}$ not contained
in $W$ which specializes to $x'$. Thus we obtain
$x \leadsto x_2 \leadsto x'$, the point $x_2$ is closed in $X_\eta$,
and $x_2 \not \in W$ as desired.
\end{proof}
\end{document}
